%%%%%%%%%%%%%%%%%%%%%%%%%%%%%%%%%%%%%%%%%%%%%%%%%%%%%%%%%%%%%%%%%%%%%%
% Lớp 11 | Hàm số lượng giác và phương trình lượng giác | [Giảm] Phương trình lượng giác thường gặp
% Dạng 1D1?6-2: Phương trình đẳng cấp bậc n đối với sinx và cosx
% Số câu: 2  (sinh tự động bởi tach_cau.py; từ nay thầy sửa trực tiếp file này)
%%%%%%%%%%%%%%%%%%%%%%%%%%%%%%%%%%%%%%%%%%%%%%%%%%%%%%%%%%%%%%%%%%%%%%

%%%%%%%%%%%%%%%% ĐÚNG/SAI %%%%%%%%%%%%%%%%

% nguồn: Ok_dungsai_lop11_Lan1
\begin{ex}%[1D1H6-2]%[ma:5c523d86]
 Gọi $x_0$ là nghiệm dương nhỏ nhất của phương trình $3\sin^2x+2\sin x\cos x-\cos^2x=0$. Chọn khẳng  định đúng?
 \choiceTF[4] 
  {\True  $x_0\in\left(0;\dfrac{\pi}{6}\right)$}
  {$x_0\in\left(\dfrac{\pi}{4};\dfrac{\pi}{2}\right)$} 
   {\True $x_0\in\left(\dfrac{\pi}{12};\dfrac{\pi}{3}\right)$}
   {$x_0\in\left(\dfrac{\pi}{6};\dfrac{\pi}{3}\right)$}
  \loigiai{
   {\allowdisplaybreaks
  \begin{eqnarray*}
   & &3\sin^2x+2\sin  x\cos x-\cos^2x=0\\
    &\Leftrightarrow&
   (3\sin x  -\cos x)(\sin x +\cos x)=0\\
    &\Leftrightarrow&
    \hoac{&3\sin x -\cos x=0\quad(1)\\&\sin x +\cos x=0.\quad(2) }
  \end{eqnarray*}
  Nhận xét $\cos x=0$ không thỏa mãn.  Khi đó
  \begin{itemize}
   \item  $(1)\Leftrightarrow\tan x=\dfrac{1}{3}\Leftrightarrow x=\arctan\dfrac{1}{3}+k\pi, (k\in\mathbb{Z})$.
   \item $(2)\Leftrightarrow\tan  x=-1\Leftrightarrow x=-\dfrac{\pi}{4}+k\pi, (k\in\mathbb{Z})$.
  \end{itemize}
  } 
  Kết hợp  hai trường hợp ta thấy nghiệm dương nhỏ nhất của phương trình đã cho là $x_0=\arctan\dfrac{1}{3}\Rightarrow x_0\in\left(\dfrac{\pi}{12};\dfrac{\pi}{6}\right)$.
  }
\end{ex}

% nguồn: Ok_dungsai_lop11_Lan1
\begin{ex}%[1D1H6-2]%[ma:9926c80f]
 Phương trình nào sau đây có tập nghiệm trùng với tập nghiệm của  phương trình $3\sin^2x+2\sin x\cos x -5\cos^2x=0$?
  \choiceTF
  {$3\cot^2x+2\cot  x-5=0$}
   {\True $5\cot^2x-2\cot x-3=0$}
  {$5\tan^2x-2\tan  x-3=0$}
   {\True $3\tan^2x+2\tan x-5=0$}
 \loigiai{\begin{itemize}
  \item $\cos x=0$, phương trình trở  thành $3=0$ (vô lý).\\
  $\cos x\neq 0$, chia hai  vế phương trình cho $\cos^2x$, phương trình đã cho trở thành $3\tan^2x+2\tan x-5=0$.
  \item  $\sin x=0$, phương trình trở thành $-5=0$ (vô lý).\\
  $\sin x\neq 0$, chia hai vế phương trình cho $\sin^2x$, phương trình đã cho trở thành $5\cot^2x-2\cot  x-3=0$.
 \end{itemize}
} 
\end{ex}
